% !TEX root = ../main.tex
\FloatBarrier\Needspace{5\baselineskip}\section{文献与证据的对应关系}

\begingroup\small\setstretch{1.06}
\begin{longtable}{@{}>{\raggedright\arraybackslash}p{0.240\dimexpr\linewidth-4\tabcolsep\relax}>{\raggedright\arraybackslash}p{0.340\dimexpr\linewidth-4\tabcolsep\relax}>{\raggedright\arraybackslash}p{0.420\dimexpr\linewidth-4\tabcolsep\relax}@{}}
\caption{文献支持与项目证据边界}\\
\toprule
文献 & 支撑内容 & 项目仍需提供的证据 \\
\midrule\endfirsthead
\toprule
文献 & 支撑内容 & 项目仍需提供的证据 \\
\midrule\endhead
\midrule\multicolumn{3}{r}{\footnotesize 续下页}\endfoot
\bottomrule\endlastfoot
\cite{ref1} Raschka；\cite{ref2} Cawley与Talbot & 评价与选择区分，小样本选择风险 & 按工件分组划分数据、固定模型及判定参数、新数据评价 \\
\cite{ref3} Saito与Rehmsmeier & PR解释及类别比例 & 实际缺陷比例及复核负担 \\
\cite{ref4} Guo等 & 温度缩放与校准动机 & 校准前后质量、融合分数可靠性 \\
\cite{ref5} Tong等；\cite{ref6} Elkan & 错误优先级、约束与成本 & 复核流程、成本和正式目标；当前无NP保证 \\
\cite{ref7} Bergstra与Bengio；\cite{ref8} Snoek等 & 后续搜索预算与昂贵实验选择 & 自身搜索过程及预算比较 \\
\cite{ref9} Shang等 & 未来生产抽检与抽样校正 & 实际抽样记录；当前区间不源于该抽样法 \\
\cite{ref10}--\cite{ref13} 模型与方法原始文献 & 骨干、多尺度与Focal机制 & 齿轮候选对照，具体参数值依据 \\
\cite{ref14} MVTec官方资源 & 工业异常检测研究设定 & 自身正常样本多样性；本项目未使用该数据集 \\
\cite{ref15} Choi与Seong；\cite{ref16} Kittler等；\cite{ref17} Wu等 & 决策融合与分类器组合的已有用法 & 本项目自身的模型输出、融合公式及验证结果；文献不是历史参数来源 \\
\end{longtable}\endgroup

\ReportAnchor{appendix-b}

\FloatBarrier\Needspace{5\baselineskip}\section{关键固定训练参数与资料索引}

\begingroup\small\setstretch{1.06}
\begin{longtable}{@{}>{\raggedright\arraybackslash}p{0.230\dimexpr\linewidth-6\tabcolsep\relax}>{\raggedright\arraybackslash}p{0.260\dimexpr\linewidth-6\tabcolsep\relax}>{\raggedright\arraybackslash}p{0.260\dimexpr\linewidth-6\tabcolsep\relax}>{\raggedright\arraybackslash}p{0.250\dimexpr\linewidth-6\tabcolsep\relax}@{}}
\caption{固定训练参数摘要}\\
\toprule
参数组 & Scratch V5 & Missing Hole V1 & 来源性质 \\
\midrule\endfirsthead
\toprule
参数组 & Scratch V5 & Missing Hole V1 & 来源性质 \\
\midrule\endhead
\midrule\multicolumn{4}{r}{\footnotesize 续下页}\endfoot
\bottomrule\endlastfoot
自定义分类初始化/输入 & ImageNet；长宽比保留、灰色填充 & 同左；输入384或512 & 外部初始化及工程输入设计 \\
归一化 & ImageNet mean/std & 同左 & 预训练输入约定 \\
自定义优化器 & AdamW，lr=3e\ensuremath{-}4，weight\_decay=5e\ensuremath{-}4 & 同左 & 脚本预先固定的参数 \\
学习率调度 & 余弦，eta\_min=1e\ensuremath{-}6 & 同左 & 脚本预先固定的参数 \\
Focal损失设置 & \ensuremath{\gamma}=1.5，最终w\_pos=1 & \ensuremath{\gamma}=1.5，w\_pos=2 & 前者候选组合，后者固定 \\
自定义batch & 384\ensuremath{\rightarrow}24；512\ensuremath{\rightarrow}12 & 同左 & 计算预算约束 \\
检测输入 & 960 & 960；1280候选未入选 & 实验变量 \\
检测优化器 & AdamW，lr0=7e\ensuremath{-}4，weight\_decay=5e\ensuremath{-}4 & 同左 & 预先固定的参数 \\
最终TTA & 无 & 无 & 验证选择 \\
\end{longtable}\endgroup

\ReportNote{表　按\protect\href{run:../GearPro_综合技术报告_Markdown含配图/GearPro_综合实验与部署技术报告.md}{《GearPro综合实验与部署技术报告》第8.3、9.4节}摘要；各实验的有效配置应结合对应args/history读取。检测结构、batch、增强与训练预算有版本差异，完整逐次表保留在综合报告。}
\ReportAnchor{project-facts}

\textbf{项目事实说明。} 本报告中由项目负责人确认的内容包括：50个齿轮工件构成、使用实际流水线工业摄像头人工摆拍多角度图像、按图像可见缺陷标注（确认不可见作为负样本，可见但难辨认或难标边界标为difficult）、仅有图片顺序编号、明确缺陷均不合格、difficult表示已确认存在但难辨认的缺陷、上下双相机、设备NX、新模型简单测试未记录、人工复核方向与流程未定、软件/模型及硬件分工、历史目标无正式需求依据。

\textbf{综合项目资料。} \protect\href{run:../GearPro_综合技术报告_Markdown含配图/GearPro_综合实验与部署技术报告.md}{GearPro综合实验与部署技术报告}，2026-10-03；其中第17节列明原始实验及阶段汇报来源。Git的main分支中的阶段PDF作为历史状态补充，当前状态以项目负责人确认的信息更新。


\begingroup\small\setstretch{1.06}
\begin{longtable}{@{}>{\raggedright\arraybackslash}p{0.220\dimexpr\linewidth-4\tabcolsep\relax}>{\raggedright\arraybackslash}p{0.330\dimexpr\linewidth-4\tabcolsep\relax}>{\raggedright\arraybackslash}p{0.450\dimexpr\linewidth-4\tabcolsep\relax}@{}}
\caption{实验资料与配置来源目录}\\
\toprule
资料名称 & 可核对文件 & 本文用途 \\
\midrule\endfirsthead
\toprule
资料名称 & 可核对文件 & 本文用途 \\
\midrule\endhead
\midrule\multicolumn{3}{r}{\footnotesize 续下页}\endfoot
\bottomrule\endlastfoot
\ReportAnchor{source-1}划痕发布配置 & \protect\href{run:../../outputs/scratch_v5/inference_config.json}{V5发布配置} & 成员、温度、融合、阈值、验证计数 \\
\ReportAnchor{source-2}齿部缺损发布配置 & \protect\href{run:../../outputs/missing_hole_v1/inference_config.json}{V1发布配置} & 同上 \\
\ReportAnchor{source-3}划痕测试记录 & \protect\href{run:../../outputs/scratch_v5/test_scratch/test_report.json}{V5测试记录} & 测试计数、相似对提示 \\
\ReportAnchor{source-4}齿部缺损及联合测试记录 & \protect\href{run:../../outputs/missing_hole_v1/test/test_report.json}{V1及联合测试记录} & 检测任务、位置、OR与统一模型 \\
\ReportAnchor{source-5}划痕数据清单 & \protect\href{run:../../dataset_defects/scratch_v5/scratch_manifest.csv}{V5数据清单} & 训练/验证拆分及缺少工件ID的状态 \\
\ReportAnchor{source-6}齿部缺损数据清单 & \protect\href{run:../../dataset_defects/missing_hole_v1/master_manifest.csv}{V1数据清单} & 各拆分、图像/标签及样例索引 \\
\ReportAnchor{source-7}划痕候选模型评价表 & \protect\href{run:../../outputs/scratch_v5/leaderboard.json}{V5候选模型评价表} & 划痕候选结果 \\
\ReportAnchor{source-8}齿部缺损候选模型评价表 & \protect\href{run:../../outputs/missing_hole_v1/leaderboard.json}{V1候选模型评价表} & 检测全因子及后续候选 \\
\ReportAnchor{source-9}齿部缺损融合候选评价表 & \protect\href{run:../../outputs/missing_hole_v1/fusion_leaderboard.json}{V1融合候选模型评价表} & 21个融合候选及最终选择 \\
\ReportAnchor{source-10}划痕逐图预测与标注 & \protect\href{run:../../outputs/scratch_v5/test_scratch/test_predictions.csv}{V5逐图预测}及原图/XML & 仅用于已记录错误的展示与解释 \\
\ReportAnchor{source-11}V3测试与定位记录 & \protect\href{run:../../outputs/defect_search_v3/final_report.json}{V3测试与定位记录} & 图像与区域两种评价、划痕计数差异 \\
\ReportAnchor{source-12}V4初始基线记录 & \protect\href{run:../../outputs/binary_defect_v4/baseline_report_before_extended.json}{V4初始基线记录} & 基线计数、指标与类别分组 \\
\ReportAnchor{source-13}V4最终融合记录 & \protect\href{run:../../outputs/binary_defect_v4/final_report.json}{V4最终融合记录} & 融合计数、指标与类别分组 \\
\ReportAnchor{source-14}V4发布配置 & \protect\href{run:../../outputs/binary_defect_v4/inference_config.json}{V4发布配置} & 成员、增强及融合方式 \\
\ReportAnchor{source-15}V4候选评价记录 & \protect\href{run:../../outputs/binary_defect_v4/leaderboard.json}{V4候选评价记录} & 扩展单模型验证阈值与增强方式 \\
\ReportAnchor{source-16}V4详细技术报告 & \protect\href{run:../V4实验详细技术报告.md}{V4详细技术报告} & 扩展单模型既有留出摘要及四路增强说明 \\
\end{longtable}\endgroup

图表由\ReportCode{build_report.py}读取上述记录生成，数值与来源汇总在\ReportCode{assets/chart_data.csv}，输入文件哈希和构建检查保存在\ReportCode{verification.json}。脚本只读取已有实验产物，不运行模型、不修改阈值。

\ReportAnchor{appendix-c}

\FloatBarrier\Needspace{5\baselineskip}\section{术语及方法名称}
正文优先使用中文名称，必要时补充英文或代码名。独立训练的网络称为“成员模型”；针对同一类缺陷的两个分类模型和一个目标检测模型合称“检测子系统”。下表解释这些名称在项目中的含义，更多代码对应见\protect\href{run:../GearPro_技术报告/术语对照表.md}{术语对照表}。

报告保留实验记录中的字段名AUPRC。实际脚本先按分数从高到低逐样本排序，再累计“召回增量\ensuremath{\times}当前精确率”，即ΣΔRecall\ensuremath{\times}Precision，属于AP式阶梯求和，并非梯形积分；相同分数按输入顺序处理，未按并列阈值成组。因此这些数值用于复现项目内的排序比较，与其他实现比较时应先统一计算方式。计算代码见\protect\href{run:../../train_scratch_v5.py}{划痕评价函数}和\protect\href{run:../../train_binary_defect_v4.py}{V4评价函数}，V1复用划痕函数。


\begingroup\small\setstretch{1.06}
\begin{longtable}{@{}>{\raggedright\arraybackslash}p{0.240\dimexpr\linewidth-4\tabcolsep\relax}>{\raggedright\arraybackslash}p{0.340\dimexpr\linewidth-4\tabcolsep\relax}>{\raggedright\arraybackslash}p{0.420\dimexpr\linewidth-4\tabcolsep\relax}@{}}
\caption{术语与实际方法定义}\\
\toprule
中文名称 & 英文或代码名 & 本项目中的定义 \\
\midrule\endfirsthead
\toprule
中文名称 & 英文或代码名 & 本项目中的定义 \\
\midrule\endhead
\midrule\multicolumn{3}{r}{\footnotesize 续下页}\endfoot
\bottomrule\endlastfoot
齿部缺损 & partial tooth damage；历史代码missing\_hole & 齿侧或齿角等部位的局部缺失；不译为孔洞缺失 \\
划痕检测子系统 & Scratch V5 & 两个分类模型与一个目标检测模型的组合 \\
齿部缺损检测子系统 & Missing Hole V1 & 两个分类模型与一个目标检测模型的组合 \\
分数级加权融合 & score-level weighted fusion & 对同类缺陷的模型分数进行加权平均 \\
决策级融合 & decision-level fusion & 对两个子系统的二值判定进行合并 \\
或规则 & OR rule & 任一子系统判为有缺陷，最终即判为有缺陷 \\
温度缩放 & temperature scaling & 用正温度参数调整模型输出分数，参数根据验证集NLL选择 \\
参考标注 & ground truth annotations & 人工标注的类别、边界框及困难目标标记 \\
困难目标 & difficult=1 & 确定存在缺陷但难以辨认的标注目标 \\
正样本/负样本 & positive/negative samples & 含/不含当前任务可见目标缺陷的图像 \\
精确率 & precision & TP/(TP+FP)，与准确率accuracy不同 \\
召回率 & recall & TP/(TP+FN)，亦可解释为缺陷检出率 \\
误报率 & false positive rate & FP/(FP+TN)，与FP/(TP+FP)不同 \\
预测为缺陷的样本比例 & predicted-positive fraction & (TP+FP)/N；图像统计不能直接换算为工件送复核比例 \\
目标级定位评价 & object-level localization evaluation & 基于预测框与标注框的匹配计算定位指标 \\
初步实验 & pilot experiment & 以较低训练轮数上限比较候选配置 \\
增加训练轮数上限的重新训练 & retraining with a larger epoch limit & 从预训练初始化重新训练，区别于从初步实验checkpoint接续训练 \\
对照实验 & comparative experiment & 比较候选配置；多项因素同时变化时不解释为单因素效应 \\
消融实验 & ablation study & 在控制其他条件下移除或改变某个组件，检验其作用；按实际设计使用 \\
\end{longtable}\endgroup

\ReportAnchor{appendix-d}

\FloatBarrier\Needspace{5\baselineskip}\section{数据检查与早期探索}
正文的数据规模见第1章；以下保留构成、标注、拆分及早期方案的详细说明。

\FloatBarrier\Needspace{4\baselineskip}\subsection{工件构成与位置标签}
项目共有50个齿轮工件：正常10个，划痕15个，侧部缺损15个，齿角缺损10个。工件由项目自行采购或定制，采用上下、旋转、侧面等多个角度拍摄。

\begin{figure}[H]\centering
\includegraphics[width=\linewidth,height=0.37\textheight,keepaspectratio]{fig03_dataset_units.png}
\caption{左侧为负责人确认的齿轮工件构成，右侧为两类缺陷的数据清单中的图像数量。多角度拍摄增加观察条件，不增加不同齿轮工件数量；图像尚未可靠对应到具体工件ID。（来源：\protect\hyperlink{project-facts}{项目事实说明}；\protect\hyperlink{source-5}{划痕数据清单}；\protect\hyperlink{source-6}{齿部缺损数据清单}）}
\label{fig:report-7}
\end{figure}

这里有三个不同概念：采购齿轮时记录的缺陷类型、拍摄时的观察角度，以及XML中的位置标签。采购分类为“侧部缺损/齿角缺损”，V1位置标注为bottom/oblique/side；其他版本还使用按宽高比推断的视角。它们之间的对应关系还没有核对，所以报告分别保留原有定义。

\FloatBarrier\Needspace{4\baselineskip}\subsection{数据规模与类别定义}

\par\medskip\noindent\begin{minipage}{\linewidth}\small\setstretch{1.06}
\captionof{table}{数据规模及各任务正负样本定义}
\begin{tabular}{@{}>{\raggedright\arraybackslash}p{0.260\dimexpr\linewidth-8\tabcolsep\relax}>{\raggedright\arraybackslash}p{0.260\dimexpr\linewidth-8\tabcolsep\relax}>{\raggedright\arraybackslash}p{0.160\dimexpr\linewidth-8\tabcolsep\relax}>{\raggedright\arraybackslash}p{0.160\dimexpr\linewidth-8\tabcolsep\relax}>{\raggedright\arraybackslash}p{0.160\dimexpr\linewidth-8\tabcolsep\relax}@{}}
\toprule
检测任务 & 训练图 & 验证图 & 测试图 & 正负样本定义 \\
\midrule
Scratch V5 & 364 & 150 & 150 & 有可见划痕 / 无可见划痕 \\
Missing Hole V1 & 397 & 99 & 150 & 有可见齿部缺损 / 无可见齿部缺损 \\
\bottomrule\end{tabular}\end{minipage}\par\medskip

\ReportNote{表　训练、验证数量来自清单（来源：\protect\hyperlink{source-5}{划痕数据清单}；\protect\hyperlink{source-6}{齿部缺损数据清单}），测试计数来自（来源：\protect\hyperlink{source-3}{划痕测试记录}；\protect\hyperlink{source-4}{齿部缺损及联合测试记录}）。两个检测子系统使用同一批150张测试图，但负样本定义不同。V1排除困难目标后的训练集另有391张图。}
150张测试图中，划痕31张、齿部缺损44张、二者同时存在4张，因此含至少一类缺陷的图像为71张，无可见划痕或齿部缺损的图像为79张。无划痕图可能具有齿部缺损，无齿部缺损图可能具有划痕；两个检测子系统FPR的分母不能直接用于联合FPR。

缺少某类标注不能自动解释为不含该类缺陷。统一模型的数据构建曾排除59张划痕状态未知且无齿部缺损的候选图像，保留类别状态已确定的数据。（来源：\protect\href{run:../GearPro_综合技术报告_Markdown含配图/GearPro_综合实验与部署技术报告.md}{《GearPro综合实验与部署技术报告》第10.1.2节}）

\FloatBarrier\Needspace{4\baselineskip}\subsection{近重复检查与工件划分限制}
已有实验用SHA256检查相同文件，并用dHash、序号、宽高比和灰度相关性检查相似图像。这些检查有助于发现重复，但当前清单只有图片顺序编号，没有照片与50个齿轮的可靠对应关系，同一齿轮的不同角度是否跨训练、验证和测试仍无法确认。（来源：\protect\hyperlink{project-facts}{项目事实说明}；\protect\hyperlink{source-5}{划痕数据清单}；\protect\hyperlink{source-6}{齿部缺损数据清单}）

部分齿轮的缺陷形态独特，可能用于辨识工件；正常齿轮和缺陷不可见的照片则无法可靠辨识。本轮没有恢复或改写工件对应关系，因此完整的工件独立性仍未知。即使可辨识的照片没有跨集重复，也还需要检查其余照片。

V5测试记录没有发现字节级重复，但提示75个dHash\ensuremath{\leq}2的跨集合相似对。齿轮外形和背景相似也可能产生这种提示，需要逐图核对；目前还不能确定是否存在数据泄漏。现有记录能够确认配置固定后进行了图像评价，按工件分组的独立性还需补查。（来源：\protect\hyperlink{source-3}{划痕测试记录}）

现有测试描述的是这些图像上的检测效果。要判断对新采购齿轮、不同批次或新生产线是否有效，后续应给每个齿轮建立ID，并将其所有角度、照明条件和连续帧放入同一个数据组。

\FloatBarrier\Needspace{4\baselineskip}\subsection{训练、验证与测试职责}
训练集用于学习网络参数，验证集用于选择模型、温度、融合方式和阈值，测试集用于评价固定配置。如果查看测试结果后继续调整方案，后续需要新的开发数据和未查看的测试数据，旧集合可用于检查修改后是否退步。\cite{ref1}\cite{ref2}

这次整理只解释已有结果，没有重新训练或调整运行阈值。各阶段的数据和标签不同，所以跨版本比较主要解释思路如何变化，性能数值仍分别报告。

\FloatBarrier\Needspace{4\baselineskip}\subsection{早期方案与转向依据}
\begin{figure}[H]\centering
\includegraphics[width=\linewidth,height=0.37\textheight,keepaspectratio]{fig04_research_path.png}
\caption{V3/V4提供研究转向依据，V5/V1承担重点技术探索。该图表示问题和设计演进，非同一数据上的性能趋势。（来源：\protect\href{run:../GearPro_综合技术报告_Markdown含配图/GearPro_综合实验与部署技术报告.md}{《GearPro综合实验与部署技术报告》第3、6--10节}）}
\label{fig:report-8}
\end{figure}

V3用分类模型区分正常、划痕和齿部缺损，另用目标检测模型定位缺陷，还尝试了增强、P2和正常特征距离。实验关注低对比划痕和不同视角下的缺损，也促使我们分别评价图像判断与区域定位。V4随后合并两类缺陷，尝试统一二分类、多模型、纹理方法和融合。（来源：\protect\href{run:../GearPro_综合技术报告_Markdown含配图/GearPro_综合实验与部署技术报告.md}{《GearPro综合实验与部署技术报告》第6--7节}）

V4最终在留出集上检出19/30张缺陷图像，误报7/20张无可见目标缺陷图像；11个漏检中7个来自划痕。这说明后续需要重点检查划痕，我们因此分别组织两类缺陷的数据与模型实验。第一次留出结果曾用于决定扩展方向，这批数据已参与探索，后续最终评价需要另选未查看的数据。（来源：\protect\href{run:../GearPro_综合技术报告_Markdown含配图/GearPro_综合实验与部署技术报告.md}{《GearPro综合实验与部署技术报告》第7.7节}）

\FloatBarrier\Needspace{4\baselineskip}\subsection{V3划痕定位计数的未解释差异}
原始测试JSON的划痕图像分类检出4/11张；整体定位评价记录TP=0、FP=6、FN=19。另一个按视角逐框匹配的结果在oblique组、IoU=0.5时记录TP=1、FP=5、FN=9，其他视角TP=0。（来源：\protect\hyperlink{source-11}{V3测试与定位记录}） 两套定位计数不能相加或互相替换，差异原因尚未核清。因此正文采用齿部缺损数据解释图像判定与定位的任务区别，不作“划痕完全无法定位”的结论。

\FloatBarrier\Needspace{4\baselineskip}\subsection{V4比较方案的配置与资料来源}

\par\medskip\noindent\begin{minipage}{\linewidth}\small\setstretch{1.06}
\captionof{table}{V4配置与历史结果来源}
\begin{tabular}{@{}>{\raggedright\arraybackslash}p{0.230\dimexpr\linewidth-6\tabcolsep\relax}>{\raggedright\arraybackslash}p{0.260\dimexpr\linewidth-6\tabcolsep\relax}>{\raggedright\arraybackslash}p{0.260\dimexpr\linewidth-6\tabcolsep\relax}>{\raggedright\arraybackslash}p{0.250\dimexpr\linewidth-6\tabcolsep\relax}@{}}
\toprule
正文名称 & 实验或配置 & 验证选择阈值（约） & 留出指标来源 \\
\midrule
基线模型 & \ReportCode{p_resnet18_512_w4}；四路增强 & 0.640567 & 基线JSON的test.overall \protect\hyperlink{source-12}{V4初始基线记录} \\
扩展单模型 & \ReportCode{xfull_resnet18_384_w1_s73_r1}；四路增强 & 0.571126 & 详细报告§7.2历史测试摘要 \protect\hyperlink{source-16}{V4详细技术报告}；阈值及增强另核对候选JSON \protect\hyperlink{source-15}{V4候选评价记录} \\
三模型融合 & 三个分类成员，分别四路增强，分数取最大值 & 0.702842 & 最终JSON的test.overall \protect\hyperlink{source-13}{V4最终融合记录}；发布配置 \protect\hyperlink{source-14}{V4发布配置} \\
\bottomrule\end{tabular}\end{minipage}\par\medskip

扩展单模型的验证AUPRC约0.860，正文表中的0.899是详细报告记载的既有留出AUPRC，不能混用；该留出摘要本轮没有逐图预测重算证据。这里保留已有记录，不运行模型或重选阈值。

基线与最终融合的划痕召回均为4/11；齿部缺损分别为13/19和15/19。两个增加的检出计数来自齿部缺损，但尚未按逐图集合确认哪些图像被补检或丢失。前向次数4与12是配置推导，不是现场计时；所有结果仍受既有留出反馈、样本规模和工件划分限制。

\ReportAnchor{appendix-e}

\FloatBarrier\Needspace{5\baselineskip}\section{候选实验与未采用配置}
\FloatBarrier\Needspace{4\baselineskip}\subsection{划痕候选与比较限制}
划痕实验尝试结合图像分类与局部检测线索。加工纹和高光可能与划痕混淆，局部预测框也可能不准，因此先比较两类模型及其分数组合对图像判断的影响，定位效果另行检查。

\textbf{对照设计。} 正式记录包含6个分类模型初步实验、2个目标检测模型初步实验和2个增加训练轮数上限的分类模型重新训练实验。初步分类实验比较ResNet18、EfficientNet-B0、YOLO分类器及384/w1、512/w2配置；目标检测实验比较960输入的原始YOLO26n与包含P2检测层的模型。实验还比较是否使用测试时增强（test-time augmentation，TTA），随后进入温度参数与融合规则选择。（来源：\protect\href{run:../GearPro_综合技术报告_Markdown含配图/GearPro_综合实验与部署技术报告.md}{《GearPro综合实验与部署技术报告》第8节}；\protect\hyperlink{source-7}{划痕候选模型评价表}）

原始YOLO26n检测模型的召回率为38/43、FPR为21/107，P2对应38/43和19/107；记录中的AUPRC约0.828与0.866（计算定义见附录C）。在相同召回率下，P2误报更少，排序指标更高，进入最终组合。最终分类成员为增加训练轮数上限后重新训练的EfficientNet-B0与初步实验中的ResNet18，两者384输入、w1、无TTA。（来源：\protect\hyperlink{source-1}{划痕发布配置}；\protect\hyperlink{source-7}{划痕候选模型评价表}）

我们根据验证结果选择P2配置，但两个检测配置的batch不同，结构的单独作用还需要对照。384/w1与512/w2也同时改变输入和权重。后续应固定其他条件比较这些设置，并核对模型在同一图像上的共同错误，检查融合成员是否必要。

P2模型的最高框置信度参与图像判断，预测框是否准确则单独评价。原报告在IoU=0.3的框匹配中记录了大量FP，所以目前显示框只供辅助观察。（来源：\protect\href{run:../GearPro_综合技术报告_Markdown含配图/GearPro_综合实验与部署技术报告.md}{《GearPro综合实验与部署技术报告》第8.8--8.9节}）

\FloatBarrier\Needspace{4\baselineskip}\subsection{齿部缺损候选与选择规则}
我们比较标签划分和困难目标处理，是为了检查有限样本下哪种设置更合适。位置标签细分后可能更容易区分外观，也可能使每类样本不足；保留或排除困难目标会改变训练标注。V1因此组合这些设置进行比较。

\textbf{对照设计。} 初步目标检测实验采用三类/两类/单类\ensuremath{\times}保留/排除困难目标\ensuremath{\times}原始YOLO26n/包含P2检测层的模型，共12组，统一960输入和35轮训练上限。随后进行12个初步分类实验、增加训练轮数上限的重新训练、第二个随机种子的验证、1280输入对照及困难样本重采样与CLAHE增强实验，总计34个实验记录。（来源：\protect\href{run:../GearPro_综合技术报告_Markdown含配图/GearPro_综合实验与部署技术报告.md}{《GearPro综合实验与部署技术报告》第9节}；\protect\hyperlink{source-8}{齿部缺损候选模型评价表}）

三类标注区分bottom/oblique/side，两类合并bottom与oblique，单类将位置归为一个齿部缺损目标。单类、排除困难目标的原始YOLO26n模型与包含P2检测层的模型均检出33/34，原始YOLO26n检测模型误报6/65，P2误报13/65。采用单类标注的原始YOLO26n模型以较低误报成为入选方向，说明细分标签或增加高分辨率特征层并未自动改善判定。

exclude\_difficult只修改训练集：删除21个困难目标的标注框，6张仅含困难目标的图像整体排除，训练图397\ensuremath{\rightarrow}391、框498\ensuremath{\rightarrow}477。这样避免把已知缺陷图误当作负样本；验证与测试保留完整参考标注。（来源：\protect\href{run:../GearPro_综合技术报告_Markdown含配图/GearPro_综合实验与部署技术报告.md}{《GearPro综合实验与部署技术报告》第9.3节}）

最终EfficientNet采用512、排除仅含困难目标的训练图像，ResNet18采用384、保留困难目标，检测器采用单类标注的原始YOLO26n模型、960、排除困难目标。不同成员保留不同训练配置；它们的互补程度仍需逐图对照。（来源：\protect\hyperlink{source-2}{齿部缺损发布配置}）

最终保留初步实验的检测权重。重新训练模型的总体召回率和FPR相同，AUPRC更高，但选择规则先比较各位置组中的最低召回率：初步实验约0.917，重新训练约0.909，所以初步模型入选。这个差别涉及少量样本，还需要重复实验检查是否稳定。（来源：\protect\href{run:../GearPro_综合技术报告_Markdown含配图/GearPro_综合实验与部署技术报告.md}{《GearPro综合实验与部署技术报告》第9.11节}）

\FloatBarrier\Needspace{4\baselineskip}\subsection{未采用配置}

\par\medskip\noindent\begin{minipage}{\linewidth}\small\setstretch{1.06}
\captionof{table}{未采用配置及已有观察}
\begin{tabular}{@{}>{\raggedright\arraybackslash}p{0.240\dimexpr\linewidth-4\tabcolsep\relax}>{\raggedright\arraybackslash}p{0.340\dimexpr\linewidth-4\tabcolsep\relax}>{\raggedright\arraybackslash}p{0.420\dimexpr\linewidth-4\tabcolsep\relax}@{}}
\toprule
探索 & 预期作用 & 现有观察与选择 \\
\midrule
V4纹理三通道 & 强调亮度、局部对比和边缘 & 已有比较未支持替换RGB输入，最终方案保留RGB \\
V5四路TTA & 减轻外观变化 & 最终模型未采用测试时增强 \\
V1 1280输入 & 增加局部细节像素 & 所比较的重新训练配置该对照未支持替换最终模型，输入规模更大 \\
V1困难样本重采样与CLAHE增强 & 改善困难区域的学习 & 未替换当前最终检测器 \\
V1增加训练轮数上限后的重新训练 & 更充分拟合数据 & 最低位置组召回率的排序未优于初步实验模型，虽AUPRC更高 \\
\bottomrule\end{tabular}\end{minipage}\par\medskip

\ReportNote{表　依据\protect\href{run:../GearPro_综合技术报告_Markdown含配图/GearPro_综合实验与部署技术报告.md}{《GearPro综合实验与部署技术报告》第7--9节}。结论针对已有配置与数据；保留这些实验，是为了说明方案怎样收敛，而非宣称相应方法在所有场景无效。}
\ReportAnchor{appendix-f}

\FloatBarrier\Needspace{5\baselineskip}\section{参数选择与补充结果}
\FloatBarrier\Needspace{4\baselineskip}\subsection{训练参数、采样与框架配置}
分类网络从ImageNet预训练参数初始化，替换输出层后训练。YOLO使用官方初始化，或从已有权重加载可迁移参数到P2配置。最终使用的checkpoint由训练保存和验证选择得到，网络各层的参数通过训练更新。（来源：\protect\href{run:../GearPro_综合技术报告_Markdown含配图/GearPro_综合实验与部署技术报告.md}{《GearPro综合实验与部署技术报告》第8--9节}）

自定义分类器使用Focal BCE：$L=\operatorname{mean}[w_y(1-p_t)^\gamma\operatorname{BCE}(z,y)]$，其中\ReportCode{p_t}为真实类别的预测分数。Focal机制降低易样本的损失占比，参考Lin等\cite{ref13}。项目预设\ensuremath{\gamma}=1.5；V5在候选组合中比较w1/w2，最终使用w1；V1固定正类权重2，其单独作用还未验证。

逆频率采样$w_i=1/n_{(\mathrm{label},\,\mathrm{view\ or\ position})}$调整样本被抽到的频率，损失权重则调整抽到后对训练损失的影响。脚本另设FN\ensuremath{\times}3、FP\ensuremath{\times}2的困难样本乘数，但V5这轮实际只对一张训练FP额外加权。因此目前还没有证据说明困难样本加权带来了主要收益。（来源：\protect\href{run:../GearPro_综合技术报告_Markdown含配图/GearPro_综合实验与部署技术报告.md}{《GearPro综合实验与部署技术报告》第8.5节}）

YOLO分类候选的w2通过复制正类文件增加正类的出现频率，自定义分类器则使用损失权重。检测args中的box/cls/dfl设置为7.5/0.5/1.5，来自框架配置，没有逐项消融；哪些字段实际参与训练，需要结合任务实现和loss日志核对。（来源：\protect\href{run:../GearPro_综合技术报告_Markdown含配图/GearPro_综合实验与部署技术报告.md}{《GearPro综合实验与部署技术报告》第5.2、8.5、9.1节}）

\FloatBarrier\Needspace{4\baselineskip}\subsection{温度参数选择}
温度缩放先将原始分数p转换为$\operatorname{logit}(p)=\ln[p/(1-p)]$，再计算$p_T=\operatorname{sigmoid}(\operatorname{logit}(p)/T)$。我们在[0.35,4.00]中比较147个等距温度候选，选择验证集负对数似然（NLL）最小的T。NLL衡量分数与真实标签的符合程度，越小越好；边界分数按发布实现处理。（来源：\protect\href{run:../GearPro_综合技术报告_Markdown含配图/GearPro_综合实验与部署技术报告.md}{《GearPro综合实验与部署技术报告》第5、8.3.4节}）

温度缩放参考Guo等\cite{ref4}。正温度不会改变单个模型分数的排序，但会改变融合时各模型分数的相对尺度。当前记录了温度选择，还需要比较校准前后的NLL、可靠性图和分箱样本量，检查分数是否更可靠。

\FloatBarrier\Needspace{4\baselineskip}\subsection{融合公式、候选与完整阈值}
将两个分类模型经温度缩放后的分数先取平均$c=(s_{\mathrm{Eff}}+s_{\mathrm{Res}})/2$，再比较$f_\alpha=\alpha c+(1-\alpha)s_{\mathrm{det}}$，\ensuremath{\alpha}\ensuremath{\in}\{0.25,0.50,0.75\}，同时比较单模型、分类分数平均/最大值及与检测模型分数取最大值的方案。V5记录20个融合候选规则，V1保留21个候选成绩。（来源：\protect\href{run:../GearPro_综合技术报告_Markdown含配图/GearPro_综合实验与部署技术报告.md}{《GearPro综合实验与部署技术报告》第8.2、9.6节}；\protect\hyperlink{source-9}{齿部缺损融合候选评价表}）

\begin{figure}[H]\centering
\includegraphics[width=\linewidth,height=0.37\textheight,keepaspectratio]{fig07_fusion_weights.png}
\caption{发布配置（来源：\protect\hyperlink{source-1}{划痕发布配置}；\protect\hyperlink{source-2}{齿部缺损发布配置}）中的融合公式展开。划痕\ensuremath{\alpha}=0.25，所以分类器分别0.125、检测器0.75；齿部缺损\ensuremath{\alpha}=0.50，所以分类器分别0.25、检测器0.50。系数表示分数计算规则，不能解释为因果贡献或定位精度。}
\label{fig:report-9}
\end{figure}


\par\medskip\noindent\begin{minipage}{\linewidth}\small\setstretch{1.06}
\captionof{table}{最终成员系数、温度及阈值}
\begin{tabular}{@{}>{\raggedright\arraybackslash}p{0.160\dimexpr\linewidth-10\tabcolsep\relax}>{\raggedright\arraybackslash}p{0.140\dimexpr\linewidth-10\tabcolsep\relax}>{\raggedright\arraybackslash}p{0.140\dimexpr\linewidth-10\tabcolsep\relax}>{\raggedright\arraybackslash}p{0.140\dimexpr\linewidth-10\tabcolsep\relax}>{\raggedright\arraybackslash}p{0.240\dimexpr\linewidth-10\tabcolsep\relax}>{\raggedright\arraybackslash}p{0.180\dimexpr\linewidth-10\tabcolsep\relax}@{}}
\toprule
检测任务 & 分类模型1：EfficientNet & 分类模型2：ResNet18 & 目标检测模型 & 温度参数T（同序） & 判定阈值（展示值） \\
\midrule
划痕V5 & 0.125 & 0.125 & 0.750 & 0.700/2.050/2.125 & 0.3003 \\
齿部缺损V1 & 0.250 & 0.250 & 0.500 & 1.000/2.375/1.850 & 0.3413 \\
\bottomrule\end{tabular}\end{minipage}\par\medskip

V1的分类平均分数与检测分数等权融合时，结果为34TP、4FP；仅将EfficientNet分数以0.25权重与检测分数组合时，计数也相同。两者AUPRC约为0.966和0.963，所以按既定排序选择前者。这是采用0.50的记录依据，其稳定优势还需要进一步比较。（来源：\protect\hyperlink{source-9}{齿部缺损融合候选评价表}）

运行时使用JSON中的完整阈值：划痕为0.300273610279458，齿部缺损为0.3413327979078584。正文显示四位小数方便阅读，运行配置保持原值。（来源：\protect\hyperlink{source-1}{划痕发布配置}；\protect\hyperlink{source-2}{齿部缺损发布配置}）

\FloatBarrier\Needspace{4\baselineskip}\subsection{阈值选择与两类结果合并}
阈值选择先扫描固定候选的验证分数，优先满足FPR\ensuremath{\leq}0.20，再提高Recall，并按规则比较位置组召回、FPR和AUPRC。齿部缺损还优先考虑各位置组中的最低召回率。这里选择的是图像判断配置，目标框的mAP需要另外评价。（来源：\protect\href{run:../GearPro_综合技术报告_Markdown含配图/GearPro_综合实验与部署技术报告.md}{《GearPro综合实验与部署技术报告》第8.3.4、9.2节}）

Recall\ensuremath{\geq}0.95、FPR\ensuremath{\leq}0.20没有正式产线需求或成本测算依据，本文将其保留为\textbf{历史实验目标及选择约束}。0.20解释当时阈值如何被选出；0.95解释历史目标，不作为本项目当前生产验收标准。结果继续保留真实数值，评价重点转向检出能力及人工复核负担。

决策级融合采用或规则：$\mathrm{REJECT}=(s_{\mathrm{scratch}}\geq\tau_{\mathrm{scratch}})\lor(s_{\mathrm{missing}}\geq\tau_{\mathrm{missing}})$。这里OR表示逻辑或：任一检测子系统判为有缺陷，最终即判为有缺陷。两个子系统的分数针对不同缺陷，负样本定义和温度参数估计所用的数据不同，因此分别进行阈值判定，再合并二值结果。每个任务的FPR约束不自动成为融合系统的FPR约束，不含这两类缺陷的图像上的误报需单独统计。工业外观检测已有采用缺陷类别独立建模与或规则决策融合的研究\cite{ref15}；该文献用于方法命名及相关工作比较，不是历史参数来源。

\FloatBarrier\Needspace{4\baselineskip}\subsection{复核评价框架}
漏检优先的使用方向需要同时呈现：\textbf{主指标召回率，配套指标误报率、精确率和预测为缺陷的样本比例}。对同一集合，预测为缺陷的样本比例为$q=(\mathrm{TP}+\mathrm{FP})/N$，它在当前记录中只表示图像预测比例。只有评价单位与工件一一对应，且全部疑似工件送复核时，才等于工件送复核比例；多视角图像不能直接作此换算。


\par\medskip\noindent\begin{minipage}{\linewidth}\small\setstretch{1.06}
\captionof{table}{人工复核评价框架}
\begin{tabular}{@{}>{\raggedright\arraybackslash}p{0.240\dimexpr\linewidth-4\tabcolsep\relax}>{\raggedright\arraybackslash}p{0.340\dimexpr\linewidth-4\tabcolsep\relax}>{\raggedright\arraybackslash}p{0.420\dimexpr\linewidth-4\tabcolsep\relax}@{}}
\toprule
评价问题 & 当前可以给出的证据 & 尚需明确的信息 \\
\midrule
有多少缺陷漏检 & FN计数、总体及困难目标所在图像的Recall & 被判为无缺陷的工件是否复核或抽查 \\
正常件增加多少检查 & FP计数和FPR & 人工复核时间、处理能力 \\
一共需要检查多少件 & 当前数据上的预测为缺陷的样本比例 & 真实缺陷比例、复核流程 \\
应选择哪个判定阈值 & 历史阈值对应的验证结果及误差分析 & 两类缺陷标签完整的新开发数据、可接受漏检与复核负担 \\
\bottomrule\end{tabular}\end{minipage}\par\medskip

成本敏感学习和错误约束研究\cite{ref5}\cite{ref6}帮助确定需要优先控制哪类错误。具体接受多少漏检、承担多少复核，还要根据实际流程确定。当前保留历史权重和阈值，后续用新的开发数据选择阈值，旧测试用于回归检查。

\FloatBarrier\Needspace{4\baselineskip}\subsection{参数来源完整表}

\begingroup\small\setstretch{1.06}
\begin{longtable}{@{}>{\raggedright\arraybackslash}p{0.240\dimexpr\linewidth-4\tabcolsep\relax}>{\raggedright\arraybackslash}p{0.340\dimexpr\linewidth-4\tabcolsep\relax}>{\raggedright\arraybackslash}p{0.420\dimexpr\linewidth-4\tabcolsep\relax}@{}}
\caption{参数来源完整表}\\
\toprule
参数 & 采用值/方法 & 来源类型与证据 \\
\midrule\endfirsthead
\toprule
参数 & 采用值/方法 & 来源类型与证据 \\
\midrule\endhead
\midrule\multicolumn{3}{r}{\footnotesize 续下页}\endfoot
\bottomrule\endlastfoot
网络参数 & 预训练初始化后微调 & 外部初始化＋对应训练实验保存 \\
分类骨干与输入 & V5两个分类模型均为384；V1为512/384 & 验证候选选择；V5输入与w同时变化 \\
检测结构 & V5含P2检测层；V1单类原始YOLO26n & 两类任务分别对照；batch等条件仍有差异 \\
Focal聚焦参数\ensuremath{\gamma} / V1正类权重 & 1.5 / 2 & 预先固定的训练参数，未独立消融 \\
采样与难例乘数 & 逆频率；FN3、FP2 & 预先设定的采样规则，实际使用范围见附录F.1 \\
温度参数 & 每个检测子系统的三个成员模型分别估计T & 验证NLL估计，数值见\protect\hyperlink{source-1}{划痕发布配置}；\protect\hyperlink{source-2}{齿部缺损发布配置} \\
融合系数 & \ensuremath{\alpha}=0.25/0.50 & 验证有限候选搜索与均值展开 \\
判定阈值 & 约0.3003/0.3413 & 固定融合的验证扫描 \\
ROI裁剪边界扩展 & 4\% & 实现设定，未独立消融 \\
候选框保留下限 / NMS IoU & 0.001 / 0.7 & 推理配置；不同于判定及框显示门槛 \\
学习率、优化器、增强范围 & 各训练脚本预先设定的值 & 完整归档见附录B，未逐项证明最优 \\
\end{longtable}\endgroup

\FloatBarrier\Needspace{4\baselineskip}\subsection{验证与测试差异}
\begin{figure}[H]\centering
\includegraphics[width=\linewidth,height=0.37\textheight,keepaspectratio]{fig09_validation_test.png}
\caption{每条线连接同一缺陷类别的验证与测试指标。两任务分母不同；验证参与多层选择，线条不表示时间上的退化。数据来自（来源：\protect\hyperlink{source-1}{划痕发布配置}；\protect\hyperlink{source-2}{齿部缺损发布配置}；\protect\hyperlink{source-3}{划痕测试记录}；\protect\hyperlink{source-4}{齿部缺损及联合测试记录}）。}
\label{fig:report-10}
\end{figure}

划痕验证集正确检出42/43、误报9/107，测试集正确检出25/31、误报20/119；齿部缺损验证集正确检出34/34、误报4/65，测试集正确检出36/44、误报4/106。测试召回均约81\%，明显低于参与选择的验证表现。

差异可能涉及验证选择偏差、拍摄条件变化和困难样本覆盖不足，当前数据不能量化各因素贡献。按图像计数的Wilson 95\%召回区间：划痕约63.7\%--90.8\%，齿部缺损约68.0\%--90.5\%；同一工件图像间的相关性未知，区间仅作描述性参考。

\FloatBarrier\Needspace{4\baselineskip}\subsection{位置分组结果}
\begin{figure}[H]\centering
\includegraphics[width=\linewidth,height=0.37\textheight,keepaspectratio]{fig11_position_recall.png}
\caption{bottom为11/12、oblique为10/17、side为20/20。位置分组可重叠，数量之和49大于44，不能相加算总体召回；区间按图像计数推导。（来源：\protect\hyperlink{source-4}{齿部缺损及联合测试记录}）}
\label{fig:report-11}
\end{figure}

齿部缺损8个漏检中7个属于oblique，说明这类位置是当前采集和可见性检查的重点。side的20/20也应结合样本量理解。包含困难目标、忽略困难目标和仅评价非困难样本这三种条件下，召回率均约0.81，困难目标标记并不能充分解释测试结果下降。（来源：\protect\href{run:../GearPro_综合技术报告_Markdown含配图/GearPro_综合实验与部署技术报告.md}{《GearPro综合实验与部署技术报告》第9.7节}；\protect\hyperlink{source-4}{齿部缺损及联合测试记录}）

\FloatBarrier\Needspace{4\baselineskip}\subsection{缺陷比例与复核负担}
联合测试可见目标缺陷图像比例71/150，是富含缺陷的研究数据，不代表生产线。假设可见缺陷图像的召回率和无可见目标缺陷图像的误报率保持不变，$\mathrm{Precision}=\pi R/[\pi R+(1-\pi)\mathrm{FPR}]$，预测为缺陷的样本比例为$q=\pi R+(1-\pi)\mathrm{FPR}$，\ensuremath{\pi}为应用条件下的可见目标缺陷图像比例，并非整件缺陷率。\cite{ref3}

测试精确率77.78\%还不足以判断产线人工复核的效果。后续应测量现场图像与工件的缺陷比例、模型未检出工件中的漏检、疑似工件的人工确认比例及复核时间，再选择适合实际流程的阈值。目前人工复核效果尚未测量。

\FloatBarrier\Needspace{4\baselineskip}\subsection{合并结果的类别解释与划痕分组统计}
齿部缺损检测本身检出36/44张正样本；两类结果合并后，这44张图像中有41张被标为疑似缺陷。其中5张由划痕检测判为有缺陷，不能据此将齿部缺损检测的正确检出数改为41。（来源：\protect\hyperlink{source-4}{齿部缺损及联合测试记录}）

历史划痕分析中，漏检组的分类模型平均分数约0.455、检测模型平均分数约0.144；误报组的分类模型平均分数约0.814。结合图像外观，加工纹和高光是值得检查的误报来源，但分组均值不能直接解释每张图像的错误原因。（来源：\protect\href{run:../GearPro_综合技术报告_Markdown含配图/GearPro_综合实验与部署技术报告.md}{《GearPro综合实验与部署技术报告》第8.6.3节}）

\FloatBarrier\Needspace{4\baselineskip}\subsection{单模型与融合对照的完整记录}
下表记录同一任务验证集上各方案按FPR\ensuremath{\leq}0.20规则选择的结果。\ensuremath{\alpha}是分类分数系数，检测分数系数为1\ensuremath{-}\ensuremath{\alpha}；若分类分数是两个分类器的均值，组合就包含三个成员。各方案使用各自阈值，组件的单独作用还需消融实验检查。划痕最终融合AUPRC未在发布配置中载明，表中保留缺项。


\begingroup\small\setstretch{1.06}
\begin{longtable}{@{}>{\raggedright\arraybackslash}p{0.070\dimexpr\linewidth-12\tabcolsep\relax}>{\raggedright\arraybackslash}p{0.200\dimexpr\linewidth-12\tabcolsep\relax}>{\raggedright\arraybackslash}p{0.170\dimexpr\linewidth-12\tabcolsep\relax}>{\raggedright\arraybackslash}p{0.070\dimexpr\linewidth-12\tabcolsep\relax}>{\raggedright\arraybackslash}p{0.240\dimexpr\linewidth-12\tabcolsep\relax}>{\raggedright\arraybackslash}p{0.110\dimexpr\linewidth-12\tabcolsep\relax}>{\raggedright\arraybackslash}p{0.140\dimexpr\linewidth-12\tabcolsep\relax}@{}}
\caption{单模型与融合完整验证记录}\\
\toprule
任务 & 实验名／配置 & 组合规则 & TP／FP & 完整阈值 & AUPRC & 来源 \\
\midrule\endfirsthead
\toprule
任务 & 实验名／配置 & 组合规则 & TP／FP & 完整阈值 & AUPRC & 来源 \\
\midrule\endhead
\midrule\multicolumn{7}{r}{\footnotesize 续下页}\endfoot
\bottomrule\endlastfoot
划痕 & \ReportCode{f_efficientnet_b0_384_w1_s20260911} & 单模型 & 41／13 & 0.542919409017341 & 0.893870 & \protect\hyperlink{source-7}{候选评价表} \\
划痕 & \ReportCode{p_detector_p2_960} & 单模型 & 38／19 & 0.18610403292287162 & 0.866130 & \protect\hyperlink{source-7}{候选评价表} \\
划痕 & 发布配置中的两个分类模型及P2检测模型 & 分类均值与检测分数加权，\ensuremath{\alpha}=0.25 & 42／9 & 0.300273610279458 & 未载明 & \protect\hyperlink{source-1}{发布配置} \\
齿部缺损 & \ReportCode{f_cls_efficientnet_b0_512_exclude_difficult_s20260913} & 单模型 & 34／6 & 0.3809993863105774 & 0.924740 & \protect\hyperlink{source-9}{融合候选评价表} \\
齿部缺损 & \ReportCode{p_det_one_class_exclude_difficult_standard_960} & 单模型 & 33／6 & 0.12059450038989711 & 0.944733 & \protect\hyperlink{source-9}{融合候选评价表} \\
齿部缺损 & \ReportCode{classifier_mean} & 两个分类分数均值 & 34／5 & 0.5251138191759659 & 0.894532 & \protect\hyperlink{source-9}{融合候选评价表} \\
齿部缺损 & \ReportCode{weighted_f_cls_efficientnet_b0_512_exclude_difficult_s20260913_0.25} & EfficientNet与检测分数加权，\ensuremath{\alpha}=0.25 & 34／4 & 0.20599650265366867 & 0.963388 & \protect\hyperlink{source-9}{融合候选评价表} \\
齿部缺损 & \ReportCode{weighted_classifier_mean_0.5} & 分类均值与检测分数加权，\ensuremath{\alpha}=0.50 & 34／4 & 0.3413327979078584 & 0.965955 & \protect\hyperlink{source-9}{融合候选评价表} \\
\end{longtable}\endgroup

本轮还未核对单模型与融合逐图的新增检出和丢失图像，因此计数差表示净变化。齿部缺损两模型与三模型都是34TP、4FP，具体预测图像、部署代价和稳定性仍需比较。当前配置保持不变，后续再检查成员必要性与NX开销。

\ReportAnchor{appendix-g}

\FloatBarrier\Needspace{5\baselineskip}\section{软件、整机状态与后续测量}
\FloatBarrier\Needspace{4\baselineskip}\subsection{软件输出链路}
我负责软件和模型，其他成员负责硬件及后续现场部署。软件分别用采集线程和推理线程处理图像，齿轮定位模型（代码名称Model1）生成ROI，再执行两个检测子系统。每次结果保存成员分数、实际阈值、缺陷类别、预测框和耗时。

\begin{figure}[H]\centering
\includegraphics[width=\linewidth,height=0.37\textheight,keepaspectratio]{fig12_software_flow.png}
\caption{软件将同次检测结果提供给页面、统计和串口；原型已完成分拣模块联调，人工复核规则和记录方式还待确定。图据\protect\href{run:../GearPro_综合技术报告_Markdown含配图/GearPro_综合实验与部署技术报告.md}{《GearPro综合实验与部署技术报告》第13节}及负责人确认绘制。}
\label{fig:report-12}
\end{figure}

软件用同一个检测结果对象（代码名称Observation）向页面、统计模块和串口提供结果，避免不同模块各自判断。加载模型时检查配置、版本和文件哈希；暂停后恢复时复用已加载模型。如果必要检测子系统运行失败或分数非法，就停止有效检测输出，避免将运行错误当成通过结果。（来源：\protect\href{run:../GearPro_综合技术报告_Markdown含配图/GearPro_综合实验与部署技术报告.md}{《GearPro综合实验与部署技术报告》第13节}）

\FloatBarrier\Needspace{4\baselineskip}\subsection{原型状态}
当前仍使用上下两台工业相机，设备已升级为NX。第一代原型已完成输送、采集、显示与分拣模块联调。新模型只做过简单现场测试，未记录测试数据，整机性能还需要系统测量。

历史软件优先处理最新帧，并设置5秒间隔抑制重复计数，以减少积压和重复统计。还需要按工件检查这些设置是否合适，尤其是一帧出现多个齿轮时，整帧判断怎样对应每件分拣。目前尚无完整跨帧工件ID追踪、检测耗时和分拣对应测量。（来源：\protect\href{run:../GearPro_综合技术报告_Markdown含配图/GearPro_综合实验与部署技术报告.md}{《GearPro综合实验与部署技术报告》第13.1--13.4节}）


\begingroup\small\setstretch{1.06}
\begin{longtable}{@{}>{\raggedright\arraybackslash}p{0.240\dimexpr\linewidth-4\tabcolsep\relax}>{\raggedright\arraybackslash}p{0.340\dimexpr\linewidth-4\tabcolsep\relax}>{\raggedright\arraybackslash}p{0.420\dimexpr\linewidth-4\tabcolsep\relax}@{}}
\caption{软件与原型状态}\\
\toprule
内容 & 当前状态 & 报告中的证据层级 \\
\midrule\endfirsthead
\toprule
内容 & 当前状态 & 报告中的证据层级 \\
\midrule\endhead
\midrule\multicolumn{3}{r}{\footnotesize 续下页}\endfoot
\bottomrule\endlastfoot
自建数据与标注 & 已完成多角度拍摄，图像与工件的对应关系未建立 & 项目负责人确认＋图像与标签清单 \\
V5/V1离线模型 & 已有配置、权重索引与测试记录 & 配置及JSON/CSV \\
Web、统计与串口 & 已实现统一结果链路 & 阶段材料记载的实现 \\
上下双相机、NX、机构 & 第一代原型链路 & 实物与负责人确认 \\
当前模型NX现场测试 & 简单测试，未记录数据 & 功能尝试，缺系统性能测量 \\
人工复核流程 & 应用方向已明确，流程未定 & 设计需求 \\
部署优化与节拍 & 尚未开展相应模型优化及时间考虑 & 后续测量与优化 \\
\end{longtable}\endgroup

\FloatBarrier\Needspace{4\baselineskip}\subsection{整机测量}
后续应从相机原帧开始记录齿轮定位、ROI、缺陷判断、串口命令和实际分拣结果。漏定位、未处理、工件错配、命令迟到和未执行都应纳入统计，并分别说明评价对象和分母。这样才能同时检查识别效果与整机流程。

NX测试先使用当前发布配置，分别记录取帧、定位、各成员模型、后处理和传输时间，同时记录设备温度、资源占用及长时运行表现，再决定是否需要FP16等优化。目前还没有确定每件工件的处理时间要求，也没有实测时延与吞吐率。

\FloatBarrier\Needspace{4\baselineskip}\subsection{后续实验与评价方法}

\par\medskip\noindent\begin{minipage}{\linewidth}\small\setstretch{1.06}
\captionof{table}{后续实验与评价方法}
\begin{tabular}{@{}>{\raggedright\arraybackslash}p{0.240\dimexpr\linewidth-4\tabcolsep\relax}>{\raggedright\arraybackslash}p{0.340\dimexpr\linewidth-4\tabcolsep\relax}>{\raggedright\arraybackslash}p{0.420\dimexpr\linewidth-4\tabcolsep\relax}@{}}
\toprule
待回答的问题 & 拟开展验证 & 主要输出 \\
\midrule
对新齿轮是否有效 & 建立工件ID；按工件和拍摄批次分组，采集训练中未见的齿轮 & 新齿轮工件测试及样本构成 \\
哪些缺陷实际拍得清 & 固定工件与相机，逐项比较照明、观察角和载台 & 原图、可见性、定位与困难子组变化 \\
V5输入与权重各自作用 & 同拆分和骨干，384/512\ensuremath{\times}w1/w2 & 解耦后的候选与波动 \\
成员模型校准和融合是否必要 & 固定成员模型参数，独立校准/开发比较，保存错误交集 & NLL、可靠性图、识别收益与开销 \\
复核流程怎样减少漏检 & 明确被判为无缺陷的工件处理、复核能力与新的判定阈值 & Recall、FPR、确认比例及复核时间 \\
NX是否满足真实节拍 & 完整原型记录分项耗时和工件命令对应 & 时延分布、覆盖、资源与稳定性 \\
\bottomrule\end{tabular}\end{minipage}\par\medskip

如果后续需要比较较多参数，可以先用随机搜索作为预算对照\cite{ref7}；训练成本较高时，再考虑贝叶斯优化\cite{ref8}。当前还没有使用这两种方法，优先把数据来源和待回答的问题确定清楚。

上线后如缺陷稀少、标注预算有限，可参考Shang等按预测结果分层抽样并校正P/R估计的研究\cite{ref9}。这涉及测试抽样，与训练正类复制不同；实际流程需包含被判为无缺陷的工件抽查，才能评价漏检。

